% Source-derived chapter 02: Algebraic side
% Original source: anticyclotomic-iwasawa-rational-elliptic-eisenstein
\chapter{Algebraic side}
\label{anticyclotomic-iwasawa-rational-elliptic-eisenstein-chapter-02}
\source{anticyclotomic-iwasawa-rational-elliptic-eisenstein}

\begin{remark}[Source section]
\label{anticyclotomic-iwasawa-rational-elliptic-eisenstein-chapter-02-source}
\source{anticyclotomic-iwasawa-rational-elliptic-eisenstein}
\source{anticyclotomic-iwasawa-rational-elliptic-eisenstein}
This chapter reproduces the corresponding section of the retrieved
LaTeX source, including its definitions, statements, examples, and
proofs. It has not yet been translated into Lean.
\notready
\end{remark}

% The source section title was: Algebraic side
\label{sec:algebraic}
\source{anticyclotomic-iwasawa-rational-elliptic-eisenstein}

In this section we prove Theorem~\ref{MAINalgside} below, relating the anticyclotomic Iwasawa invariants of an elliptic curve $E/\Q$ at a prime $p$ with $E[p]^{ss}=\mathbb{F}_p(\phi)\oplus\mathbb{F}(\psi)$ to the anticyclotomic Iwasawa invariants of the characters $\phi$ and $\psi$. 

Throughout, we fix a prime $p>2$ and an embedding $\iota_p:\overline{\Q}\hookrightarrow\overline{\Q}_p$, and let $K\subset\overline{\Q}$ be an imaginary quadratic field in which $p=v\bar{v}$ splits, with $v$ the prime of $K$ above $p$ induced by $\iota_p$.  We also fix an embedding $\iota_\infty:\overline{\Q}\hookrightarrow \C$.

Let $G_K = {\rm Gal}(\overline{\Q}/K)\subset G_\Q = {\rm Gal}(\overline{\Q}/\Q)$, and for each place $w$ of $K$ let $I_w\subset G_w\subset G_K$ be
corresponding inertia and decomposition groups. Let $\Frob_w\in G_w/I_w$ be the arithmetic Frobenius. For the prime $v\mid p$ we assume 
$G_v$ is chosen so that it is identified with ${\rm Gal}(\overline{\Q}_p/\Q_p)$ via $\iota_p$.

%We also 
Let $\Gamma={\rm Gal}(K_\infty/K)$ be the Galois group of the anticyclotomic $\Z_p$-extension $K_\infty$ of $K$, and 
let $\Lambda=\Z_p\llbracket\Gamma\rrbracket$ be the anticyclotomic Iwasawa algebra. We shall often identify $\Lambda$ with the power series ring $\bZ_p\llbracket T\rrbracket$ by setting $T=\gamma-1$ for a fixed topological generator $\gamma\in\Gamma$.


\subsection{Local cohomology groups of characters}\label{localchar}
\source{anticyclotomic-iwasawa-rational-elliptic-eisenstein}


Let  $\theta:G_K\rightarrow\mathbb{F}_p^\times$  
be a character with conductor divisible only by primes that are split in $K$. Via the Teichm\"uller lift $\mathbb{F}_p^\times\hookrightarrow\Z_p^\times$, we shall also view $\theta$ as taking values in $\Z_p^\times$. Set
\[
M_\theta=\Z_p(\theta)\otimes_{\Z_p}\Lambda^\vee,
\]
where $(-)^\vee = {\rm Hom}_{\rm cts}(-,\bQ_p/\bZ_p)$ for topological $\Z_p$-modules. The module $M_\theta$ is equipped with a $G_K$-action via $\theta\otimes\Psi^{-1}$, where $\Psi:G_K\rightarrow\Lambda^\times$ is the character arising from the projection $G_K\twoheadrightarrow\Gamma$.

In this section, we study the local cohomology of $M_\theta$ at various primes $w$ of $K$. 

\subsubsection{$w\nmid p$ split in $K$}

Let $w$ be a prime of $K$ lying over a prime $\ell\neq p$ split in $K$, and let $\Gamma_w\subset \Gamma$ be the corresponding decomposition group. Let $\gamma_w\in\Gamma_w$ be 
the image of $\Frob_w$, and set
\begin{equation}\label{eq:Euler}
\source{anticyclotomic-iwasawa-rational-elliptic-eisenstein}
\mathcal{P}_w(\theta)=P_w(\ell^{-1}\gamma_w)
%=(1-\theta(\ell)\ell^{-1}\gamma_w)
\in\Lambda,
\end{equation}
where $P_w={\rm det}(1-{\rm Frob}_wX\vert\bQ_p(\theta)_{I_w})$ is the Euler factor at $w$ of the $L$-function of $\theta$.
 
\begin{lemma}
\source{anticyclotomic-iwasawa-rational-elliptic-eisenstein}
\notready\label{lemmawneqp} 
\source{anticyclotomic-iwasawa-rational-elliptic-eisenstein}
The module $\rH^1(K_w,M_\theta)^\vee$ is $\Lambda$-torsion with
\[
{\rm char}_\Lambda(\rH^1(K_w,M_\theta)^\vee)=(\mathcal{P}_w(\theta)).
\]
In particular, $\rH^1(K_w,M_\theta)^\vee$ has $\mu$-invariant zero.
\end{lemma}

\begin{proof} 
Since $\ell$ splits in $K$, it follows from class field theory that the index $[\Gamma:\Gamma_w]$ is finite (i.e., $w$ is finitely decomposed in $K_\infty/K$). Thus the argument proving \cite[Prop.~2.4]{greenvats} can be immediately adapted to yield this result.
\end{proof}


\subsubsection{$w\mid p$} Recall that we assume that $p=v\bar{v}$ splits in $K$. We begin by recording the following commutative algebra lemma, which shall also be used later in the paper.

\begin{lemma}
\source{anticyclotomic-iwasawa-rational-elliptic-eisenstein}
\notready\label{lem:commalg}
\source{anticyclotomic-iwasawa-rational-elliptic-eisenstein}
Let $X$ be a finitely generated $\Lambda$-module satisfying the following two properties: 
\begin{itemize}
\item $X[T]=0$, 
\item $X/TX$ is a free $\Z_p$-module of rank $r$. 
\end{itemize}
Then $X$ is a free $\Lambda$-module of rank $r$.
\end{lemma}

\begin{proof}
From Nakayama's lemma we obtain a surjection $\pi:\Lambda^r\twoheadrightarrow X$ which becomes an isomorphism $\bar{\pi}$ after reduction modulo $T$. Letting $K={\rm ker}(\pi)$, from the snake lemma we deduce the exact sequence
\[
0\rightarrow K/TK\rightarrow(\Lambda/T\Lambda)^r\xrightarrow{\bar{\pi}}X/TX\rightarrow 0.
\]
Thus $K/TK=0$, and so $K=0$ by another application of Nakayama's lemma. 
\end{proof}

Let $\omega:G_\Q\rightarrow \mathbb{F}_p^\times$ be the mod $p$ cyclotomic character.  
Let $w$ be a prime of $K$ above $p$.

\begin{prop}
\source{anticyclotomic-iwasawa-rational-elliptic-eisenstein}
\notready\label{propw=v}
\source{anticyclotomic-iwasawa-rational-elliptic-eisenstein}
Assume that $\theta\vert_{G_w}\neq\mathds{1},\omega$. Then:
\begin{enumerate}
\item[(i)] The restriction map
\[
r_w:\rH^1(K_w,M_\theta)\rightarrow\rH^1(I_w,M_\theta)^{G_w/I_w}
\]
is an isomorphism. %In particular, $\rH^1_{\rm unr}(K_v,M_E)=0$.
\item[(ii)] $\rH^1(K_w,M_{\theta})$ is $\Lambda$-cofree of rank $1$.
\end{enumerate}
\end{prop}

\begin{proof} 
The map $r_w$ is clearly surjective, so it suffices to show injectivity. Since $G_w/I_w$ is pro-cyclic,
\[
{\rm ker}(r_w)\simeq M_\theta^{I_w}/({\rm Frob}_w-1)M_\theta^{I_w},
\]
where ${\rm Frob}_w$ is a Frobenius element at $w$. Taking Pontryagin duals to the exact sequence
\[
0\rightarrow M_\theta^{G_w}\rightarrow M_\theta^{I_w}\xrightarrow{{\rm Frob}_w-1}M_\theta^{I_w}\rightarrow M_\theta^{I_w}/({\rm Frob}_w-1)M_\theta^{I_w}\rightarrow 0
\]
and using the vanishing of $M_\theta^{G_w}$ (which follows from $\theta\vert_{G_w}\neq\mathds{1}$) we deduce a $\Lambda$-module surjection 
\begin{equation}\label{eq:surj}
\source{anticyclotomic-iwasawa-rational-elliptic-eisenstein}
(M_\theta^\vee)_{I_w}\twoheadrightarrow(M_\theta^\vee)_{I_w},
\end{equation}
hence an isomorphism (by the Noetherian property of $\Lambda$).
Since the kernel of (\ref{eq:surj}) is isomorphic to ${\rm ker}(r_w)^\vee$, (i)  follows. For (ii), in light of Lemma~\ref{lem:commalg}, letting 
\[
X:=\rH^1(K_w,M_\theta)^\vee,
\] 
it suffices to show that $X[T]=0$ and the quotient $X/TX$ is $\bZ_p$-free of rank $1$. Taking cohomology for the exact sequence $0\rightarrow\bQ_p/\bZ_p(\theta)\rightarrow M_{\theta}\xrightarrow{\times T}M_{\theta}\rightarrow 0$ we obtain 
\begin{equation}\label{eq:coh-char}
\source{anticyclotomic-iwasawa-rational-elliptic-eisenstein}
\frac{\rH^1(K_w,M_{\theta})}{T\rH^1(K_w,M_{\theta})}=0,\quad
\quad\rH^1(K_w,\Q_p/\Z_p(\theta))\simeq\rH^1(K_w,M_{\theta})[T],
\end{equation}
using that $\rH^2(K_w,M_\theta)=0$ (which follows from $\theta|_{G_w} \neq \omega$)
for the first isomorphism and 
$\rH^0(K_w,M_{\theta})=0$ for the second. The first isomorphism shows that $X[T]=0$. On the other hand, taking cohomology for the exact sequence $0\rightarrow\mathbb{F}_p(\theta)\rightarrow\bQ_p/\bZ_p(\theta)\xrightarrow{p}\bQ_p/\bZ_p(\theta)\rightarrow 0$ and using that $\theta\vert_{G_w}\neq\omega$ we obtain
\[
\frac{\rH^1(K_w,\bQ_p/\bZ_p(\theta))}{p\rH^1(K_w,\bQ_p/\bZ_p(\theta))}\simeq\rH^2(K_w,\mathbb{F}_p(\theta))=0,
\]
which together with the second isomorphism in (\ref{eq:coh-char}) shows that $X/TX\simeq\rH^1(K_w,\bQ_p/\bZ_p(\theta))^\vee$ is $\bZ_p$-free of rank $1$ (the value of the rank following from the local Euler characteristic formula), concluding the proof.
\end{proof}


\subsection{Selmer groups of characters}
\label{subsec:Sel-char}
\source{anticyclotomic-iwasawa-rational-elliptic-eisenstein}

As in the preceding section, let $\theta:G_K\rightarrow\mathbb{F}_p^\times$ be a character whose conductor is divisible only by primes split in $K$ (that is, which are unramified over $\Q$ and have degree one).

Let $\Sigma$ be a finite set of places of $K$ containing $\infty$ and the primes dividing $p$ or the conductor of $\theta$
and such that every finite place in $\Sigma$ is split in $K$, and denote by $K^\Sigma$ the maximal extension of $K$ unramified outside $\Sigma$. 

\begin{definition}
\source{anticyclotomic-iwasawa-rational-elliptic-eisenstein}
\notready 
The \emph{Selmer group} of $\theta$ is
\[
\rH^1_{\Fcal_{\rm Gr}}(K,M_\theta):=\ker\biggr\{\rH^1(K^\Sigma/K,M_\theta)\rightarrow\prod_{w \in \Sigma, w \nmid p}\rH^1(K_w,M_\theta)\times\rH^1(K_{\bar{v}},M_\theta)\biggr\},
\]
and letting $S=\Sigma\setminus\{v,\bar{v},\infty\}$, we define the \emph{$S$-imprimitive Selmer group} of $\theta$ by 
\[
\rH^1_{\Fcal_{\rm Gr}^S}(K,M_\theta):=\ker\biggr\{\rH^1(K^\Sigma/K,M_\theta)\rightarrow\rH^1(K_{\bar{v}},M_\theta)\biggr\}.
\]
Replacing $M_\theta$ by $M_\theta[p]$ in the above definitions, we obtain the \emph{residual Selmer group} $\rH^1_{\Fcal_{\rm Gr}}(K,M_\theta[p])$ and its $S$-imprimitive variant $\rH^1_{\Fcal_{\rm Gr}^S}(K,M_\theta[p])$.
\end{definition}

It is well-known that the above  groups are cofinitely generated over the corresponding Iwasawa algebra ($\Lambda$ and $\Lambda/p$), and that the Selmer group and residual Selmer groups 
are independent of the choice of the set $\Sigma$ as above. 

%Our definition of these Selmer groups is slightly different from those considered \cite{greenvats} (where the unramified, rather than strict, local condition is imposed), and is more convenient for our applications. Nonetheless, as we exploit in the next result, one can show that both definitions give rise to the same Selmer groups under mild hypotheses on $\theta$ by building on the local results in the preceding section.

The following result, combining work of Rubin  
and Hida, will play a key role in our proofs.

\begin{theorem}[Rubin, Hida]
\source{anticyclotomic-iwasawa-rational-elliptic-eisenstein}
\notready\label{thmmu=0} 
\source{anticyclotomic-iwasawa-rational-elliptic-eisenstein}
Assume that $\theta\vert_{G_{\bar{v}}}\neq\mathds{1},\omega$. Then $\rH^1_{\Fcal_{\rm Gr}}(K,M_\theta)^\vee$ is a torsion $\Lambda$-module with $\mu$-invariant zero.
\end{theorem}

\begin{proof}
Let $K_\theta\subset\overline{\bQ}$ be the fixed field of ${\rm ker}(\theta)$, and set $\Delta_\theta={\rm Gal}(K_\theta/K)$. The restriction map
\[
\rH^1(K^\Sigma/K,M_\theta)\rightarrow\rH^1(K^\Sigma/K_\theta,M_\theta)^{\Delta_\theta}
\] 
is an isomorphism (since $p\nmid\vert\Delta_\theta\vert$), which combined with Shapiro's lemma gives rise to an identification
\begin{equation}\label{eq:CFT}
\source{anticyclotomic-iwasawa-rational-elliptic-eisenstein}
\rH^1(K^\Sigma/K_\theta,M_\theta)
%\simeq{\rm Hom}_{\Delta_\theta}(K^\Sigma/K,\bQ_p/\bZ_p)
\simeq{\rm Hom}_{\rm cts}((\mathcal{X}_\infty^\Sigma)^\theta,\bQ_p/\bZ_p),
\end{equation}
where $\mathcal{X}_\infty^\Sigma={\rm Gal}(\mathcal{M}_\infty^\Sigma/K_\infty K_\theta)$ is the Galois group of the maximal abelian pro-$p$ extension of $K_\infty K_\theta$ unramified outside $\Sigma$, and $(\mathcal{X}_\infty^\Sigma)^\theta$ is the $\theta$-isotypic component of $\mathcal{X}_\infty^\Sigma$ for the action of $\Delta_\theta$, identified as a subgroup of ${\rm Gal}(K_\infty K_\theta/K)$ via the decomposition ${\rm Gal}(K_\infty K_\theta/K)\simeq\Gamma\times\Delta_\theta$. 

Now, by \cite[Rem.~3.2]{pollack2011anticyclotomic} (since the primes $w\nmid p$ in $\Sigma$ are finitely decomposed in $K_\infty/K$) and Proposition~\ref{propw=v}(i), the Selmer group $\rH^1_{\Fcal_{\rm Gr}}(K,M_\theta)$ is the same as the one defined by the unramified local conditions, i.e., as
\[
\ker\biggr\{\rH^1(K^\Sigma/K,M_\theta)\rightarrow\prod_{w \in \Sigma, w \nmid p}\rH^1(I_w,M_\theta)^{G_v/I_v}\times\rH^1(I_{\bar{v}},M_\theta)^{G_{\bar{v}}/I_{\bar{v}}}\biggr\},
\]
and so under the identification $(\ref{eq:CFT})$ we obtain
\[
\rH^1_{\Fcal_{\rm Gr}}(K,M_\theta)\simeq{\rm Hom}_{\rm cts}(\mathcal{X}_\infty^\theta,\bQ_p/\bZ_p)
\] 
where $\mathcal{X}_\infty={\rm Gal}(\mathcal{M}_\infty/K_\infty K_\theta)$ is the Galois group of the maximal abelian pro-$p$ extension of $K_\infty K_\theta$ unramified outside $v$. Thus from the works of  Rubin \cite{rubinmainconj},
which identifies $\mathrm{char}_\Lambda(\rH^1_{\Fcal_{\rm Gr}}(K,M_\theta)^\vee)$ with the ideal generated by an anticyclotomic projection of a Katz $p$-adic $L$-function,  
and Hida \cite{hidamu=0}, proving the vanishing of the $\mu$-invariant of such anticyclotomic $p$-adic $L$-functions, we obtain the theorem.
\end{proof}

\begin{remark}
\source{anticyclotomic-iwasawa-rational-elliptic-eisenstein}
\notready
Following the notations introduced in the proof of Theorem~\ref{thmmu=0}, and letting $\mathcal{X}_\infty^{sp}={\rm Gal}(\mathcal{M}_\infty^{sp}/K_\infty K_\theta)$ be the Galois group of the maximal abelian pro-$p$ extension of $K_\infty K_\theta$ unramified outside $v$ and in which the primes above $\bar{v}$ split completely, Proposition~\ref{propw=v}(i) shows $(\mathcal{X}_\infty^\Sigma)^\theta=(\mathcal{X}_\infty^{sp})^\theta$. 
\end{remark}

The next two results will allow us to determine $\lambda(\mathfrak{X}_\theta^S)$ in terms of the residual Selmer group $\rH^1_{\Fcal_{\rm Gr}^S}(K,M_\theta[p])$. In brief, the fact that $\mathfrak{X}_\theta^S$ has no nonzero pseudo-null $\Lambda$-submodules (shown in Proposition~\ref{propchar} below) yields the equality $\lambda(\mathfrak{X}_\theta^S)=\mathrm{dim}_{\mathbb{F}_p}\bigl(\rH^{1}_{\Fcal_{\rm Gr}^S}(K,M_\theta)[p]\bigr)$, which combined with the next lemma yields the desired result.

\begin{lemma}
\source{anticyclotomic-iwasawa-rational-elliptic-eisenstein}
\notready\label{rmkchar}
\source{anticyclotomic-iwasawa-rational-elliptic-eisenstein}
Assume that $\theta\vert_{G_{\bar{v}}}\neq\mathds{1}$. Then
\[
\rH^1_{\Fcal_{\rm Gr}^S}(K,M_\theta[p])\simeq
\rH^1_{\Fcal_{\rm Gr}^S}(K,M_\theta)[p].
\] 
\end{lemma}

\begin{proof}
The hypothesis on $\theta$ implies in particular that $\rH^0(K,\mathbb{F}_p(\theta))=0$, and so $\rH^0(K,M_\theta)=0$. Thus the natural map
\[
\rH^1(K^\Sigma/K,M_\theta[p])\rightarrow\rH^1(K^\Sigma/K,M_\theta)[p]
\]
induced by multiplication by $p$ on $M_\theta$ is an isomorphism. To conclude it suffices to check that the natural map $r_{\bar{v}}:\rH^1(K_{\bar{v}},M_\theta[p])\rightarrow\rH^1(K_{\bar{v}},M_\theta)[p]$ is an injection, but since $\rH^0(K_{\bar{v}},\mathbb{F}_p(\theta))=0$ by the hypothesis, the same argument as above shows that $r_{\bar{v}}$ is an isomorphism.
\end{proof}

Let
\[
\mathfrak{X}_\theta^S:=\rH^1_{\Fcal_{\rm Gr}^S}(K,M_\theta)^\vee \ \ \text{and} \ \ 
%,\quad
\mathfrak{X}_\theta:=\rH^1_{\Fcal_{\rm Gr}}(K,M_\theta)^\vee,
\]
and recall the element $\mathcal{P}_w(\theta)\in\Lambda$ introduced in $(\ref{eq:Euler})$.

\begin{prop}
\source{anticyclotomic-iwasawa-rational-elliptic-eisenstein}
\notready\label{propchar} 
\source{anticyclotomic-iwasawa-rational-elliptic-eisenstein}
Assume that $\theta\vert_{G_{\bar{v}}}\neq\mathds{1},\omega$. Then
$\mathfrak{X}_{\theta}^{S}$ is a torsion $\Lambda$-module with $\mu$-invariant zero and 
its $\lambda$-invariant satisfies
\begin{displaymath}
\lambda\bigl(\mathfrak{X}_{\theta}^{S}\bigr)=\lambda\bigl(\mathfrak{X}_{\theta}\bigr)+\sum_{w\in\Sigma, w\nmid p}\lambda\bigl(\mathcal{P}_w(\theta)\bigr).
\end{displaymath} 
Moreover, $\rH^{1}_{\Fcal^{S}_{\rm Gr}}(K,M_\theta[p])$ is finite and 
\[
\mathrm{dim}_{\mathbb{F}_p}\bigl(\rH^{1}_{\Fcal_{\rm Gr}^S}(K,M_\theta[p])\bigr)=\lambda\bigl(\mathfrak{X}_\theta^S\bigr).
\]
\end{prop}

\begin{proof}
Since $\mathfrak{X}_{\theta}$ is $\Lambda$-torsion by Theorem~\ref{thmmu=0} and the Cartier dual ${\rm Hom}(\bQ_p/\bZ_p(\theta),\mu_{p^\infty})$ has no non-trivial $G_{K_\infty}$-invariants, from \cite[Prop.~A.2]{pollack2011anticyclotomic} we obtain that the restriction map in the definition of $\rH^{1}_{\Fcal_{\rm Gr}}(K,M_\theta)$ is surjective, and so the sequence
\begin{equation}\label{eq:sur1}
\source{anticyclotomic-iwasawa-rational-elliptic-eisenstein}
 0 \rightarrow \rH^{1}_{\Fcal_{\rm Gr}}(K,M_\theta) \rightarrow\rH^{1}(K^{\Sigma}/K, M_\theta) \rightarrow \prod_{w\in\Sigma, w\nmid p}\rH^1(K_w,M_\theta) \times \rH^{1}(K_{\bar{v}}, M_\theta)\rightarrow 0
\end{equation}
is exact. From the definitions, this readily yields the exact sequence 
\begin{equation}\label{eq:sur2}
\source{anticyclotomic-iwasawa-rational-elliptic-eisenstein}
0 \rightarrow\rH^{1}_{\Fcal_{\rm Gr}}(K,M_\theta) \rightarrow\rH^{1}_{\Fcal_{\rm Gr}^S}(K,M_\theta) \rightarrow \prod_{w\in S}\rH^1(K_w,M_\theta) \rightarrow 0,
\end{equation}
which combined with Theorem~\ref{thmmu=0} and Lemma~\ref{lemmawneqp}  gives the first part of the proposition.

For the second part, note that $\rH^{2}(K^{\Sigma}/K,M_\theta)=0$. (Indeed, by the Euler characteristic formula, the $\Lambda$-cotorsionness of $\rH^{1}_{\Fcal_{\rm Gr}}(K,M_\theta)$ implies that $\rH^{2}(K^{\Sigma}/K,M_\theta)$ is $\Lambda$-cotorsion; being $\Lambda$-cofree, as follows immediately from the fact that ${\rm Gal}(K^\Sigma/K)$ has cohomological dimension $2$, it must vanish.) Thus from the long exact sequence in cohomology induced by 
$0\rightarrow \bQ_p/\bZ_p(\theta)\rightarrow M_\theta\xrightarrow{\times T}M_\theta\rightarrow 0$ 
we obtain the isomorphism
\[
\frac{\rH^{1}(K^{\Sigma}/K,M_\theta)}{T\rH^{1}(K^{\Sigma}/K,M_\theta)} \simeq\rH^{2}(K^{\Sigma}/K,\bQ_p/\bZ_p(\theta)).
\]

Since $\rH^{2}(K^{\Sigma}/K,\bQ_p/\bZ_p(\theta))$ is $\bZ_p$-cofree (because ${\rm Gal}(K^\Sigma/K)$ has cohomological dimension $2$), it follows that $\rH^{1}(K^{\Sigma}/K,M_\theta)^\vee$ has no nonzero pseudo-null $\Lambda$-submodules (\emph{cf.} \cite[Prop.~5]{greenberg-iwasawa}), and since $(\ref{eq:sur1})$ and $(\ref{eq:sur2})$ readily imply that
\[
\mathfrak{X}_\theta^{S}\simeq\frac{\rH^1(K^\Sigma/K,M_\theta)^\vee}{\rH^1(K_{\bar{v}},M_\theta)^\vee}
\]
as $\Lambda$-modules, by Proposition~\ref{propw=v}(iii) and \cite[Lem.~2.6]{greenvats} we conclude that also $\mathfrak{X}_\theta^S$ has no nonzero pseudo-null $\Lambda$-submodules. Finally, since $\mathfrak{X}_\theta^S$ is $\Lambda$-torsion with $\mu$-invariant zero by Theorem~\ref{thmmu=0}, the finiteness of $\rH^{1}_{\Fcal^{S}_{\rm Gr}}(K,M_\theta)[p]$ (and therefore of $\rH^{1}_{\Fcal^{S}_{\rm Gr}}(K,M_\theta[p])$ by Lemma~\ref{rmkchar}) follows from the structure theorem. It also follows that 
%$\mathfrak{X}_\theta^{S}$ is a finitely generated $\bZ_p$-module, and since $\mathfrak{X}_\theta^{S}$ has no nonzero pseudo-null $\Lambda$-submodules, we conclude that $\mathfrak{X}_\theta^S$ is a free $\bZ_p$-modules and hence that 
$\rH^1_{\mathcal{F}_{\rm Gr}^S}(K,M_\theta)$ is divisible. In particular,
\[
\rH^1_{\mathcal{F}_{\rm Gr}^S}(K,M_\theta)\simeq(\bQ_p/\bZ_p)^{\lambda},%\quad\quad\textrm{where $\lambda=\lambda(X_\theta^S)$,}
\]
where $\lambda=\lambda(\mathfrak{X}_\theta^S)$, which together with Lemma~\ref{rmkchar} gives the final formula for the $\lambda$-invariant.
\end{proof}

The following corollary will be used crucially in the next section.

\begin{cor}
\source{anticyclotomic-iwasawa-rational-elliptic-eisenstein}
\notready\label{corcharacters} 
\source{anticyclotomic-iwasawa-rational-elliptic-eisenstein}
Assume that $\theta\vert_{G_{\bar{v}}}\neq\mathds{1},\omega$. Then 
$\rH^{2}(K^{\Sigma}/K,M_\theta[p])=0$ and the sequence
\[
0\rightarrow\rH^1_{\Fcal_{\rm Gr}^S}(K,M_\theta[p])\rightarrow\rH^1(K^\Sigma/K,M_\theta[p])\rightarrow\rH^1(K_{\bar{v}},M_\theta[p])\rightarrow 0
\]
is exact.
\end{cor}

\begin{proof} 
In the course of the proof of Proposition~\ref{propchar} we showed that $\rH^2(K^\Sigma/K,M_\theta)=0$, and so the cohomology long exact sequence induced by multiplication by $p$ on $M_\theta$ yields an isomorphism
\begin{equation}\label{eq:iso}
\source{anticyclotomic-iwasawa-rational-elliptic-eisenstein}
\frac{\rH^1(K^\Sigma/K,M_\theta)}{p\rH^1(K^\Sigma/K,M_\theta)}\simeq\rH^2(K^\Sigma/K,M_\theta[p]).
\end{equation}
On the other hand, from the exactness of $(\ref{eq:sur1})$  we deduce the exact sequence
\begin{equation}\label{eq:sur3}
\source{anticyclotomic-iwasawa-rational-elliptic-eisenstein}
0\rightarrow\rH^1_{\Fcal_{\rm Gr}^S}(K,M_\theta)\rightarrow\rH^1(K^\Sigma/K,M_\theta)\rightarrow\rH^1(K_{\bar{v}},M_\theta)\rightarrow 0.
\end{equation}
Since we also showed in that proof that $\rH^1_{\Fcal_{\rm Gr}^S}(K,M_\theta)$ is divisible, and $\rH^1(K_{\bar{v}},M_\theta)$ is $\Lambda$-cofree by Proposition~\ref{propw=v}(ii), it follows from $(\ref{eq:sur3})$ that $\rH^1(K^\Sigma/K,M_\theta)^\vee$ has no $p$-torsion, and so 
\[
\rH^2(K^\Sigma/K,M_\theta[p])=0
\] 
by $(\ref{eq:iso})$, giving the first claim in the statement.

For the second claim, consider the commutative diagram
\begin{center}
\begin{tikzcd}[column sep=2em, row sep=2em]
0 \arrow[r] & \rH^1_{\Fcal_{\rm Gr}^S}(K,M_\theta) \arrow[d, "p"] \arrow[r]  \arrow[r] & \rH^{1}(K^{\Sigma}/K, M_\theta) \arrow[r] \arrow[d, "p"]  \arrow[r] & \rH^{1}(K_{\bar{v}},M_\theta) \arrow[d, "p"] \arrow[r]  & 0 \\
0 \arrow[r] & \rH^1_{\Fcal_{\rm Gr}^S}(K,M_\theta) \arrow[r] \arrow[r] & \rH^{1}(K^\Sigma/K,M_\theta) \arrow[r]  & \rH^{1}(K_{\bar{v}},M_\theta) \arrow[r] & 0,
\end{tikzcd}
\end{center}
where the vertical maps are the natural ones induced by multiplication by $p$ on $M_\theta$. Since $\rH^1_{\Fcal_{\rm Gr}^S}(K,M_\theta)$ is divisible, the snake lemma applied to this diagram yields the exact sequence
\[
0\rightarrow\rH^1_{\Fcal_{\rm Gr}^S}(K,M_\theta)[p]\rightarrow\rH^1(K^\Sigma/K,M_\theta)[p]\rightarrow\rH^1(K_{\bar{v}},M_\theta)[p]\rightarrow 0,
\]
which by Lemma~\ref{rmkchar} (and the natural isomorphisms shown in its proof) is identified with the exact sequence in the statement.
\end{proof}


\subsection{Local cohomology groups of \texorpdfstring{$E$}{E}}
\label{subsec:local-E}
\source{anticyclotomic-iwasawa-rational-elliptic-eisenstein}

Now we let $E/\mathbb{Q}$ be an elliptic curve of conductor $N$ with good reduction at $p$ and admitting a rational $p$-isogeny. %Thus 
The $G_\bQ$-module $E[p]$ is therefore reducible, fitting into an exact sequence 
\begin{equation}\label{eq:es-mod}
\source{anticyclotomic-iwasawa-rational-elliptic-eisenstein}
0\rightarrow\mathbb{F}_p(\phi)\rightarrow E[p]\rightarrow\mathbb{F}_p(\psi)\rightarrow 0,
\end{equation}
where $\phi,\psi:G_\bQ\rightarrow\mathbb{F}_p^\times$ are characters such that $\phi\psi=\omega$ by the Weil pairing. We assume that every prime $\ell\vert N$ splits in $K$ and continue to assume that $p=v\bar{v}$ splits in $K$, so the results of the preceding sections can be applied to the restrictions of $\phi$ and $\psi$ to $G_K$.

Let $T=T_pE$ be the $p$-adic Tate module of $E$, and denote by $M_E$ the $G_K$-module
\[
M_E:=T\otimes_{\bZ_p}\Lambda^\vee,
\]
where the tensor product is endowed with the diagonal $G_K$-action (and the action on $\Lambda^\vee$ is via $\Psi^{-1}$, as before). 


\begin{lemma}
\source{anticyclotomic-iwasawa-rational-elliptic-eisenstein}
\notready\label{lemma411}
\source{anticyclotomic-iwasawa-rational-elliptic-eisenstein}
Let $w$ be a prime of $K$ above $p$, and assume that $E(K_w)[p]=0$. Then $\rH^1(K_w,M_E)$ is $\Lambda$-cofree of rank $2$.
\end{lemma}

\begin{proof} 
The proof is virtually the same as the proof of Proposition~\ref{propw=v}(ii). Letting $X:=\rH^1(K_w,M_E)^\vee$, by Lemma~\ref{lem:commalg} it suffices to show that $X[T]=0$ and $X/TX$ is $\bZ_p$-free of rank $2$. The hypotheses imply that $E(K_w)[p^\infty]=0$, and so $\rH^2(K_w,E[p^\infty])=0$ by local duality. Taking cohomology for the exact sequence 
\[
0\rightarrow E[p^\infty]\rightarrow M_E\xrightarrow{\times T}M_E\rightarrow 0
\] 
it follows that
\begin{equation}\label{eq:coh}
\source{anticyclotomic-iwasawa-rational-elliptic-eisenstein}
\frac{\rH^1(K_w,M_E)}{T\rH^1(K_w,M_E)}\simeq\rH^2(K_w,E[p^\infty])=0,\quad
\quad\rH^1(K_w,E[p^\infty])\simeq\rH^1(K_w,M_E)[T].
\end{equation}

The first isomorphism shows that $X[T]=0$. On the other hand, taking cohomology for the exact sequence $0\rightarrow E[p]\rightarrow E[p^\infty]\xrightarrow{p}E[p^\infty]\rightarrow 0$ we obtain
\[
\frac{\rH^1(K_w,E[p^\infty])}{p\rH^1(K_w,E[p^\infty])}\simeq\rH^2(K_w,E[p])=0,
\]
which together with the second isomorphism in (\ref{eq:coh}) shows that $X/TX\simeq\rH^1(K_w,E[p^\infty])^\vee$ is $\bZ_p$-free. That its rank is $2$ follows from the local Euler characteristic formula.
\end{proof}


\subsection{Selmer groups of \texorpdfstring{$E$}{E}}
\label{subsec:Sel-E}
\source{anticyclotomic-iwasawa-rational-elliptic-eisenstein}

Fix a finite set $\Sigma$ of places of $K$ containing $\infty$ and the primes above $Np$, and such that the finite places in $\Sigma$ are all split in $K$. 

Similarly as in $\S\ref{subsec:Sel-char}$, we define a Selmer group for $E$ by  
\[
\rH^1_{\Fcal_{\rm Gr}}(K,M_E):=\ker\biggr\{\rH^1(K^\Sigma/K,M_E)\rightarrow\prod_{w \in \Sigma, w \nmid p}\rH^1(K_w,M_E)\times\rH^1(K_{\bar{v}},M_E)\biggr\},
\]
and an 
{$S$-imprimitive Selmer group}, where $S=\Sigma\setminus\{v,\bar{v},\infty\}$, by 
\[
\rH^1_{\Fcal_{\rm Gr}^S}(K,M_E):=\ker\biggr\{\rH^1(K^\Sigma/K,M_E)\rightarrow\rH^1(K_{\bar{v}},M_E)\biggr\}.
\]
The residual Selmer groups $\rH^1_{\Fcal_{\rm Gr}}(K,M_E[p])$ and $\rH^1_{\Fcal_{\rm Gr}^S}(K,M_E[p])$ are defined in the same manner.

Viewing the characters $\phi$ and $\psi$ appearing in the exact sequence $(\ref{eq:es-mod})$ as taking values in $\Z_p^\times$ via the Teichm\"uller lift, we obtain an exact sequence 
\begin{equation}\label{eq:esIw}
\source{anticyclotomic-iwasawa-rational-elliptic-eisenstein}
0 \rightarrow M_\phi[p] \rightarrow M_E[p] \rightarrow M_\psi[p] \rightarrow 0
\end{equation}
of $\mathrm{Gal}(K^{\Sigma}/K)$-modules. Let $G_p\subset G_\Q$ be a decomposition group at $p$.

\begin{prop}
\source{anticyclotomic-iwasawa-rational-elliptic-eisenstein}
\notready\label{cormodp} 
\source{anticyclotomic-iwasawa-rational-elliptic-eisenstein}
Assume that $\phi\vert_{G_{p}}\neq\mathds{1},\omega$. Then 
$(\ref{eq:esIw})$ induces a natural exact sequence 
\[
0 \rightarrow \rH^1_{\Fcal_{\rm Gr}^S}(K,M_\phi[p]) \rightarrow \rH^1_{\Fcal_{\rm Gr}^S}(K,M_E[p]) \rightarrow 
\rH^1_{\Fcal_{\rm Gr}^S}(K,M_\psi[p]) \rightarrow 0.
\]
In particular, $\rH^1_{\Fcal_{\rm Gr}^S}(K,M_E[p])$ is finite, and 
\[
\mathrm{dim}_{\mathbb{F}_p}\bigl(\rH^1_{\Fcal_{\rm Gr}^S}(K,M_E[p])\bigr)=\lambda(\mathfrak{X}_\phi^S)+\lambda(\mathfrak{X}_{\psi}^S).
\]
\end{prop}

\begin{proof} 
Taking cohomology for the exact sequence $(\ref{eq:esIw})$ we obtain the commutative diagram 
\begin{center}
\begin{tikzcd}[column sep=2em, row sep=2em]
0 \arrow[r] & \rH^{1}(K^{\Sigma}/K,M_\phi[p]) \arrow[d] \arrow[r] & \rH^{1}(K^{\Sigma}/K, M_E[p]) \arrow[r] \arrow[d]  & \rH^{1}(K^{\Sigma}/K,M_\psi[p]) \arrow[d]  \arrow[r] & 0 \\
0 \arrow[r] & \rH^{1}(K_{\bar{v}},M_\phi[p]) \arrow[r] & \rH^{1}(K_{\bar{v}},M_E[p]) \arrow[r]  & \rH^{1}(K_{\bar{v}},M_\psi[p])
\arrow[r] & 0,
\end{tikzcd}
\end{center}
where the exactness of the rows follows immediately from Corollary~\ref{corcharacters} and the hypothesis on $\phi$
(which implies that $\psi|_{G_p}\neq\mathds{1},\omega$ as well), 
and the vertical maps are given by restriction. Since the left vertical arrow is surjective by Corollary~\ref{corcharacters}, the snake lemma applied to this diagram yields the exact sequence in the statement. The last claim now follows from the last claim of Proposition~\ref{propchar}.
\end{proof}

Now we can relate the imprimitive residual and $p^\infty$-Selmer groups. Set
\[ 
\mathfrak{X}_{E}^S:=\rH^1_{\Fcal_{\rm Gr}^S}(K,M_E)^{\vee}, \quad \mathfrak{X}_{E}:=\rH^1_{\Fcal_{\rm Gr}}(K,M_E)^{\vee}.
\]

\begin{prop}
\source{anticyclotomic-iwasawa-rational-elliptic-eisenstein}
\notready\label{propmodp}
\source{anticyclotomic-iwasawa-rational-elliptic-eisenstein}
Assume that $\phi\vert_{G_{p}}\neq\mathds{1},\omega$. Then 
\[
\rH^1_{\Fcal_{\rm Gr}^S}(K,M_E[p]) \simeq \rH^1_{\Fcal_{\rm Gr}^S}(K,M_E)[p]. 
\]
Moreover, the modules $\mathfrak{X}_E^S$ and $\mathfrak{X}_E$ are both $\Lambda$-torsion with $\mu=0$.
\end{prop}

\begin{proof} 
Since $\psi=\omega\phi^{-1}$, our assumption on $\phi$ implies that $E(K_{\bar{v}})[p]=0$, and therefore $\rH^0(K_{\bar{v}},M_E)=0$. Thus the same argument as in the proof of Lemma~\ref{rmkchar} yields the  isomorphism in the statement. It follows from Proposition~\ref{cormodp} that $\rH^1_{\Fcal_{\rm Gr}^S}(K,M_E)[p]$ is finite, and so $\mathfrak{X}_E^S$ is $\Lambda$-cotorsion with $\mu=0$. Since $\mathfrak{X}_E$ is a quotient of $\mathfrak{X}_E^S$, this completes the proof.
\end{proof}

Now we can deduce the following analogue of Proposition~\ref{propchar} for $M_E$.

\begin{cor}
\source{anticyclotomic-iwasawa-rational-elliptic-eisenstein}
\notready\label{prop426} 
\source{anticyclotomic-iwasawa-rational-elliptic-eisenstein}
Assume that $\phi\vert_{G_{p}}\neq\mathds{1},\omega$. Then $\mathfrak{X}_E^S$ has no non-trivial finite $\Lambda$-submodules, and 
\[
\lambda\bigl(\mathfrak{X}_E^S\bigr)=\mathrm{dim}_{\mathbb{F}_p}\bigl(\rH^1_{\Fcal_{\rm Gr}^S}(K,M_E[p])\bigr).
\] 
\end{cor}

\begin{proof} 
Since $M_E^*={\rm Hom}(M_E,\mu_{p^\infty})$ has no non-trivial $G_K$-invariants and $\mathfrak{X}_E^S$ is $\Lambda$-torsion by Proposition~\ref{propmodp}, from \cite[Prop.~A.2]{pollack2011anticyclotomic} we deduce that the sequence
\begin{equation}\label{eq:1}
\source{anticyclotomic-iwasawa-rational-elliptic-eisenstein}
0 \rightarrow \rH^1_{\Fcal_{\rm Gr}}(K,M_E) \rightarrow \rH^{1}(K^{\Sigma}/K, M_E) \rightarrow \prod_{w\in\Sigma, w\nmid p}\rH^1(K_w,M_E)\times\rH^{1}(K_{\bar{v}}, M_E) \rightarrow 0
\end{equation}
is exact. Proceeding as in the proof of Proposition~\ref{propchar}, we see that the $\Lambda$-torsionness of $\mathfrak{X}_E$ implies that $\rH^{2}(K^{\Sigma}/K, M_E)=0$ and that $\rH^{1}(K^{\Sigma}/K, M_E)^{\vee}$ has no nonzero pseudo-null $\Lambda$-submodules. The exactness of $(\ref{eq:1})$ readily implies a $\Lambda$-module isomorphism 
\[
\mathfrak{X}_E^S\simeq\frac{\rH^1(K^\Sigma/K,M_E)^\vee}{\rH^{1}(K_{\bar{v}}, M_E)^\vee}.
\]
 
Since $\rH^1(K_{\bar{v}},M_E)$ is $\Lambda$-cofree by Lemma~\ref{lemma411}, we thus conclude from \cite[Lem.~2.6]{greenvats} that $\mathfrak{X}_E^S$ has no nonzero finite $\Lambda$-submodules. Together with the isomorphism  $\rH^1_{\Fcal_{\rm Gr}^S}(K,M_E[p]) \simeq \rH^1_{\Fcal_{\rm Gr}^S}(K,M_E)[p]$ of Proposition~\ref{propmodp}, the last claim in the statement
of the corollary follows from this.
\end{proof}


Finally, we note that as in Lemma~\ref{lemmawneqp}, one can show that for primes $w\nmid p$ split in $K$, the module $\rH^1(K_w,M_E)^\vee$ is $\Lambda$-torsion with characteristic ideal generated by the element
\[
\mathcal{P}_w(E)=P_w(\ell^{-1}\gamma_w)
%=(1-\theta(\ell)\ell^{-1}\gamma_w)
\in\Lambda,
\]
where $P_w={\rm det}(1-{\rm Frob}_wX\vert V_{I_w})$, for $V=T\otimes\bQ_p$, is the Euler factor at $w$ of the $L$-function of $E$. 


\subsection{Comparison I: Algebraic Iwasawa invariants}

We now arrive at the main result of this section. Recall that every prime $w\in\Sigma\setminus\{\infty\}$ is split in $K$, and we set $S=\Sigma\setminus\{v,\bar{v},\infty\}$. 

\begin{theorem}
\source{anticyclotomic-iwasawa-rational-elliptic-eisenstein}
\notready\label{MAINalgside}
\source{anticyclotomic-iwasawa-rational-elliptic-eisenstein}
Assume that $\phi\vert_{G_p}\neq\mathds{1},\omega$. Then the module $\mathfrak{X}_{E}$ is $\Lambda$-torsion with $\mu(\mathfrak{X}_{E})=0$ and 
\begin{displaymath}
\lambda\bigl(\mathfrak{X}_E\bigr)=\lambda\bigl(\mathfrak{X}_\phi\bigr)+\lambda\bigl(\mathfrak{X}_\psi\bigr)
+\sum_{w\in S} \lbrace \lambda\bigl(\mathcal{P}_w(\phi)\bigr)+\lambda\bigl(\mathcal{P}_w(\psi)\bigr)-\lambda\bigl(\mathcal{P}_w(E)\bigr) \rbrace.
\end{displaymath}
\end{theorem}

\begin{proof}
That $\mathfrak{X}_E$ is $\Lambda$-torsion with $\mu$-invariant zero is part of Proposition~\ref{propmodp}. For the $\lambda$-invariant, combining Corollary~\ref{prop426} and the last claim of Proposition~\ref{cormodp} we obtain
\begin{equation}\label{eq:lambda-imp}
\source{anticyclotomic-iwasawa-rational-elliptic-eisenstein}
\lambda\bigl(\mathfrak{X}_E^S\bigr)
%=\mathrm{dim}_{\mathbb{F}_p}\bigl(\rH^1_{\Fcal_{\rm Gr}^S}(K,M_E[p])\bigr)
=\lambda\bigl(\mathfrak{X}_\phi^S\bigr)+\lambda\bigl(\mathfrak{X}_{\psi}^S\bigr).
\end{equation}
On the other hand, from (\ref{eq:1}) we deduce the exact sequence
\[
0 \rightarrow \rH^1_{\Fcal_{\rm Gr}}(K,M_E) \rightarrow \rH^1_{\Fcal_{\rm Gr}^S}(K,M_E) \rightarrow \prod_{w\in S}\rH^1(K_w,M_E) \rightarrow 0, 
\]
and therefore the relation $\lambda\bigl(\mathfrak{X}_E^S\bigr)=\lambda\bigl(\mathfrak{X}_E\bigr)+\sum_{w\in S}\lambda\bigl(\mathcal{P}_w(E)\bigr)$. This, combined with the second part of Proposition~\ref{propchar} shows that $(\ref{eq:lambda-imp})$ reduces to the equality of $\lambda$-invariants in the statement
of the theorem.
\end{proof}
